\documentclass[12pt]{article}
\usepackage{amsmath}
\usepackage{amsfonts}
\usepackage{amssymb}
\usepackage{physics}
\usepackage{geometry}
\geometry{a4paper, margin=0.7in}


\title{02YMECH - Solution 2}
\date{Assigned for the week of Sep 29, 2025}
\author{}

\begin{document}
\maketitle

\section*{Solutions}

\begin{enumerate}

% 1
\item \textbf{Uniform circular motion.}

A particle moves counter-clockwise on a circle of radius $R = 5~\text{m}$ with speed $v = 10~\text{m/s}$. The center of the circle is at $(x_c, y_c) = (3,4)~\text{m}$. At $t=0$, the particle is at $(8,4)$, i.e. $5$ units to the right of the center, on the positive $x$-side.

\begin{enumerate}
    \item \textit{Position vector as a function of time.}

    Angular speed:
    \[
    \omega = \frac{v}{R} = \frac{10}{5} = 2~\text{rad/s}.
    \]

    Because the motion is counter-clockwise and starts at angle $0$ relative to the $+x$ direction from the center, we can parametrize:
    \[
    x(t) = x_c + R \cos(\omega t) = 3 + 5\cos(2t),
    \]
    \[
    y(t) = y_c + R \sin(\omega t) = 4 + 5\sin(2t).
    \]

    Therefore, the position vector is
    \[
    \mathbf{r}(t)
    = \big(3 + 5\cos(2t)\big)\,\hat{\imath}
    + \big(4 + 5\sin(2t)\big)\,\hat{\jmath}.
    \]

    At $t=0$,
    \[
    \mathbf{r}(0)
    = (3+5,\,4+0)
    = (8,4),
    \]
    which matches the given initial condition.

    \item \textit{Velocity, acceleration at $t=\pi/2$, and perpendicularity.}

    Velocity is time derivative:
    \begin{align*}
    \mathbf{v}(t)
    &= \dv{\mathbf{r}}{t} \\
    &= \dv{t}\left[(3 + 5\cos(2t))\,\hat{\imath}
    + (4 + 5\sin(2t))\,\hat{\jmath}\right] \\
    &= (-10\sin(2t))\,\hat{\imath}
    + (10\cos(2t))\,\hat{\jmath}.
    \end{align*}

    Its magnitude:
    \[
    \|\mathbf{v}(t)\|
    = \sqrt{(-10\sin(2t))^2 + (10\cos(2t))^2}
    = \sqrt{100(\sin^2(2t)+\cos^2(2t))}
    = 10~\text{m/s},
    \]
    as expected.

    Acceleration is the derivative of velocity:
    \begin{align*}
    \mathbf{a}(t)
    &= \dv{\mathbf{v}}{t} \\
    &= \dv{t}\left[(-10\sin(2t))\,\hat{\imath}
    + (10\cos(2t))\,\hat{\jmath}\right] \\
    &= (-20\cos(2t))\,\hat{\imath}
    + (-20\sin(2t))\,\hat{\jmath}.
    \end{align*}

    Evaluate at $t = \pi/2$:
    \[
    2t = 2\cdot \frac{\pi}{2} = \pi,
    \quad
    \sin(\pi)=0,
    \quad
    \cos(\pi)=-1.
    \]

    Then
    \[
    \mathbf{v}\left(\frac{\pi}{2}\right)
    = (-10\sin\pi)\,\hat{\imath} + (10\cos\pi)\,\hat{\jmath}
    = (0)\,\hat{\imath} + (10 \cdot (-1))\,\hat{\jmath}
    = -10\,\hat{\jmath}~\text{m/s},
    \]
    \[
    \mathbf{a}\left(\frac{\pi}{2}\right)
    = (-20\cos\pi)\,\hat{\imath} + (-20\sin\pi)\,\hat{\jmath}
    = (-20\cdot (-1))\,\hat{\imath} + (-20 \cdot 0)\,\hat{\jmath}
    = 20\,\hat{\imath}~\text{m/s}^2.
    \]

    Now check the dot product at $t=\pi/2$:
    \[
    \mathbf{v}\left(\frac{\pi}{2}\right)\cdot
    \mathbf{a}\left(\frac{\pi}{2}\right)
    = (-10\,\hat{\jmath})\cdot(20\,\hat{\imath})
    = 0.
    \]
    So they are perpendicular at that instant.

    In fact, for \emph{all} $t$,
    \begin{align*}
    \mathbf{v}(t)\cdot \mathbf{a}(t)
    &= (-10\sin(2t))(-20\cos(2t))
    + (10\cos(2t))(-20\sin(2t)) \\
    &= 200 \sin(2t)\cos(2t)
    -200 \sin(2t)\cos(2t) \\
    &= 0.
    \end{align*}
    Thus $\mathbf{v}(t)\perp \mathbf{a}(t)$ for all time. This is characteristic of uniform circular motion: acceleration is purely centripetal (radially inward), velocity is purely tangential.
\end{enumerate}


% 2
\item \textbf{Two particles in the plane.}

The position vectors are
\[
\mathbf{r}_1(t) = (1+2t)\,\hat{\imath} + (3-t)\,\hat{\jmath}, 
\qquad
\mathbf{r}_2(t) = (4-t)\,\hat{\imath} + (1+3t)\,\hat{\jmath}.
\]

\begin{enumerate}
    \item \textit{Velocities and accelerations.}

    Differentiate each with respect to $t$.

    For particle 1:
    \[
    \mathbf{v}_1(t)
    = \dv{\mathbf{r}_1}{t}
    = 2\,\hat{\imath} + (-1)\,\hat{\jmath}
    = 2\,\hat{\imath} - \hat{\jmath},
    \]
    \[
    \mathbf{a}_1(t)
    = \dv{\mathbf{v}_1}{t}
    = \mathbf{0}.
    \]
    So particle 1 moves at constant velocity, no acceleration.

    For particle 2:
    \[
    \mathbf{v}_2(t)
    = \dv{\mathbf{r}_2}{t}
    = (-1)\,\hat{\imath} + 3\,\hat{\jmath}
    = -\hat{\imath} + 3\,\hat{\jmath},
    \]
    \[
    \mathbf{a}_2(t)
    = \dv{\mathbf{v}_2}{t}
    = \mathbf{0}.
    \]
    So particle 2 also moves at constant velocity, no acceleration.

    \item \textit{Time of closest approach and minimum distance.}

    The displacement from particle 2 to particle 1 is
    \[
    \Delta \mathbf{r}(t)
    = \mathbf{r}_1(t) - \mathbf{r}_2(t).
    \]
    Compute components:
    \begin{align*}
    \Delta x(t)
    &= (1+2t) - (4-t)
    = 1+2t-4+t
    = 3t - 3,
    \\
    \Delta y(t)
    &= (3-t) - (1+3t)
    = 3-t-1-3t
    = 2 - 4t.
    \end{align*}

    Thus
    \[
    \Delta \mathbf{r}(t)
    = (3t-3)\,\hat{\imath} + (2-4t)\,\hat{\jmath}.
    \]

    The square of the distance between them is
    \[
    D^2(t)
    = \|\Delta \mathbf{r}(t)\|^2
    = (3t-3)^2 + (2-4t)^2.
    \]

    Expand:
    \begin{align*}
    (3t-3)^2 &= 9t^2 -18t +9, \\
    (2-4t)^2 &= 16t^2 -16t +4.
    \end{align*}

    So
    \[
    D^2(t)
    = (9t^2 -18t +9) + (16t^2 -16t +4)
    = 25t^2 -34t +13.
    \]

    To minimize the distance, we minimize $D^2(t)$ (same $t$ since square root is monotonic). Take derivative:
    \[
    \dv{D^2}{t} = 50t - 34.
    \]
    Set to zero:
    \[
    50t - 34 = 0
    \quad\Rightarrow\quad
    t = \frac{34}{50} = \frac{17}{25} = 0.68~\text{s}.
    \]

    Now evaluate the minimum distance. First compute $D^2$ at $t = 17/25$:
    \[
    D^2\!\left(\frac{17}{25}\right)
    = 25\left(\frac{17}{25}\right)^2 
    - 34\left(\frac{17}{25}\right)
    + 13.
    \]

    Compute term by term:
    \[
    \left(\frac{17}{25}\right)^2 = \frac{289}{625},
    \quad
    25 \cdot \frac{289}{625} = \frac{7225}{625} = 11.56.
    \]
    Next:
    \[
    34 \cdot \frac{17}{25} = \frac{578}{25} = 23.12.
    \]
    So
    \[
    D^2\!\left(\frac{17}{25}\right)
    = 11.56 - 23.12 + 13
    = 1.44.
    \]

    Therefore
    \[
    D_\text{min} = \sqrt{1.44} = 1.2~\text{m}.
    \]

    Summary: the two particles are closest at
    \[
    t = \frac{17}{25}~\text{s} = 0.68~\text{s},
    \qquad
    D_\text{min} = 1.2~\text{m}.
    \]

    Interpretation: since both have constant (but different) velocities, their relative motion is uniform, so the closest approach happens when the line between them is perpendicular to the relative velocity.
\end{enumerate}


% 3
\item \textbf{Explosion and conservation of momentum.}

Total mass before breakup:
\[
M = 13~\text{kg}.
\]
Masses after breakup:
\[
m_1 = 3~\text{kg}, \quad m_2 = 4~\text{kg}, \quad m_3 = 6~\text{kg},
\]
with speeds immediately after breakup:
\[
\mathbf{v}_1 = (4~\text{m/s})\,\hat{\imath} \quad \text{(east)},
\qquad
\mathbf{v}_2 = (3~\text{m/s})\,\hat{\jmath} \quad \text{(north)}.
\]

Before breakup, the body was at rest, so initial total momentum was zero:
\[
\mathbf{p}_\text{total, initial} = \mathbf{0}.
\]
Momentum is conserved, so the vector sum of the fragment momenta after breakup must still be zero:
\[
m_1 \mathbf{v}_1 + m_2 \mathbf{v}_2 + m_3 \mathbf{v}_3 = \mathbf{0}.
\]
Solve for $\mathbf{v}_3$:
\[
m_3 \mathbf{v}_3 = - \left( m_1 \mathbf{v}_1 + m_2 \mathbf{v}_2 \right).
\]

First compute the known momenta:
\[
m_1 \mathbf{v}_1 = (3~\text{kg})(4~\text{m/s})\,\hat{\imath}
= 12\,\hat{\imath}~\text{kg}\cdot\text{m/s},
\]
\[
m_2 \mathbf{v}_2 = (4~\text{kg})(3~\text{m/s})\,\hat{\jmath}
= 12\,\hat{\jmath}~\text{kg}\cdot\text{m/s}.
\]

Thus
\[
m_1 \mathbf{v}_1 + m_2 \mathbf{v}_2
= 12\,\hat{\imath} + 12\,\hat{\jmath}
\quad (\text{kg}\cdot\text{m/s}).
\]

So
\[
m_3 \mathbf{v}_3
= - (12\,\hat{\imath} + 12\,\hat{\jmath})
= (-12)\,\hat{\imath} + (-12)\,\hat{\jmath}.
\]

Since $m_3=6~\text{kg}$,
\[
\mathbf{v}_3
= \frac{1}{6}(-12\,\hat{\imath} - 12\,\hat{\jmath})
= (-2)\,\hat{\imath} + (-2)\,\hat{\jmath}
\quad \text{m/s}.
\]

So the third fragment moves with velocity
\[
\mathbf{v}_3 = (-2, -2)~\text{m/s}.
\]

Speed:
\[
v_3 = \sqrt{(-2)^2 + (-2)^2}
= \sqrt{4+4}
= \sqrt{8}
= 2\sqrt{2} \approx 2.83~\text{m/s}.
\]

Direction: both components are negative, so it moves toward the southwest, at $45^\circ$ below the $-x$ axis (i.e. equally south and west).

We can express the direction angle $\theta$ measured from due west toward south:
\[
\tan \theta = \frac{2}{2} = 1
\quad\Rightarrow\quad
\theta = 45^\circ.
\]

Thus $m_3$ moves to the southwest at $2\sqrt{2}~\text{m/s}$.


% 4
\item \textbf{Block on horizontal surface with time-dependent force.}

Given:
\[
m = 4~\text{kg}, \qquad
\mu_s = 0.5, \qquad
\mu_k = 0.3, \qquad
g=10~\text{m/s}^2.
\]
Applied horizontal force:
\[
F(t) = 2t \quad \text{N}.
\]

Normal force:
\[
N = mg = (4)(10) = 40~\text{N}.
\]

Maximum static friction magnitude:
\[
f_s^{\max} = \mu_s N = 0.5 \times 40 = 20~\text{N}.
\]

Kinetic friction magnitude (once sliding):
\[
f_k = \mu_k N = 0.3 \times 40 = 12~\text{N}.
\]

\begin{enumerate}
    \item \textit{When does the block start to move?}

    While the block is at rest, friction adjusts to match the applied force $F(t)$ up to a maximum of $20~\text{N}$. The block will start moving when $F(t)$ exceeds $f_s^{\max}$.

    We solve
    \[
    F(t) = 20~\text{N}
    \quad\Rightarrow\quad
    2t = 20
    \quad\Rightarrow\quad
    t = 10~\text{s}.
    \]

    Interpretation: for $0 \le t < 10~\text{s}$, static friction is enough to hold it in place, so velocity is zero. At $t=10~\text{s}$, motion begins.

    \item \textit{Acceleration immediately after it starts moving.}

    Just after $t = 10~\text{s}$, the block is in motion, so kinetic friction applies with constant magnitude $f_k = 12~\text{N}$ opposing the motion.

    The applied force at $t=10~\text{s}$ is
    \[
    F(10) = 2(10) = 20~\text{N}.
    \]

    Net horizontal force just after motion begins:
    \[
    F_{\text{net}} = F - f_k = 20 - 12 = 8~\text{N}.
    \]

    Newton's second law:
    \[
    a = \frac{F_{\text{net}}}{m}
    = \frac{8}{4}
    = 2~\text{m/s}^2.
    \]

    So the block's initial acceleration (once it breaks free) is $2~\text{m/s}^2$.

    \item \textit{Velocity 5 seconds after it starts moving.}

    After $t=10~\text{s}$, the block continues to be pushed by $F(t)=2t$. For $t>10~\text{s}$, it is sliding, so friction is kinetic and constant at $12~\text{N}$.

    For $t>10~\text{s}$:
    \[
    F_{\text{net}}(t) = F(t) - f_k = 2t - 12.
    \]
    Then the acceleration for $t>10$ is
    \[
    a(t) = \frac{F_{\text{net}}(t)}{m}
    = \frac{2t - 12}{4}
    = \frac{1}{2}t - 3
    \quad \text{(m/s$^2$)}.
    \]

    
    Integrate from $10$ to $15$ (the velocity at $t=10$ is $0$):
    \begin{align*}
    v(15) - v(10)
    &= \int_{10}^{15} \left( \frac{1}{2}t - 3 \right) dt \\
    &= \left[ \frac{1}{4}t^2 -3t \right]_{10}^{15} \\
    &= \frac{65}{4}~\text{m/s}.
    \end{align*}
\end{enumerate}

\end{enumerate}


\end{document}
